\documentclass[11pt,a4paper]{article}

% --- Encoding and fonts ---
\usepackage[utf8]{inputenc}
\usepackage[T1]{fontenc}

% --- Mathematics ---
\usepackage{amsmath,amssymb,amsthm}
\usepackage{mathtools}

% --- Page geometry ---
\usepackage[margin=1in,headheight=14pt]{geometry}

% --- Hyperlinks ---
\usepackage[colorlinks=true,linkcolor=red,citecolor=blue,urlcolor=blue]{hyperref}

% --- Lists ---
\usepackage{enumitem}

% --- Colors and boxes ---
\usepackage{xcolor}
\usepackage[theorems,skins,breakable]{tcolorbox}

% --- Header/Footer ---
\usepackage{fancyhdr}
\pagestyle{fancy}
\fancyhf{}
\fancyhead[L]{\textsc{Lesson 8: Convex Analysis}}
\fancyhead[R]{\thepage}
\renewcommand{\headrulewidth}{0.4pt}

% --- Tables ---
\usepackage{booktabs}

% --------------------------------------------------------------------------
% Theorem environments
% --------------------------------------------------------------------------
\newtheorem{theorem}{Theorem}[section]
\newtheorem{lemma}[theorem]{Lemma}
\newtheorem{proposition}[theorem]{Proposition}
\newtheorem{corollary}[theorem]{Corollary}

\theoremstyle{definition}
\newtheorem{definition}[theorem]{Definition}
\newtheorem{example}[theorem]{Example}
\newtheorem{exercise}[theorem]{Exercise}

\theoremstyle{remark}
\newtheorem{remark}[theorem]{Remark}

% --------------------------------------------------------------------------
% Colored boxes
% --------------------------------------------------------------------------
\newtcolorbox{keyresult}[1][]{%
  colback=green!5!white,
  colframe=green!50!black,
  fonttitle=\bfseries,
  title={Key Result},
  breakable,
  #1
}

\newtcolorbox{rlconnection}[1][]{%
  colback=blue!5!white,
  colframe=blue!50!black,
  fonttitle=\bfseries,
  title={RL/ML Connection},
  breakable,
  #1
}

\newtcolorbox{warning}[1][]{%
  colback=red!5!white,
  colframe=red!50!black,
  fonttitle=\bfseries,
  title={Warning},
  breakable,
  #1
}

% --------------------------------------------------------------------------
% Custom commands
% --------------------------------------------------------------------------
\newcommand{\R}{\mathbb{R}}
\newcommand{\N}{\mathbb{N}}
\newcommand{\Q}{\mathbb{Q}}
\newcommand{\Z}{\mathbb{Z}}
\newcommand{\C}{\mathbb{C}}
\newcommand{\Prob}{\mathbb{P}}
\newcommand{\E}{\mathbb{E}}
\newcommand{\indicator}{\mathbf{1}}

\DeclareMathOperator{\dom}{dom}
\DeclareMathOperator{\epi}{epi}
\DeclareMathOperator{\conv}{conv}
\DeclareMathOperator{\relint}{relint}
\DeclareMathOperator{\aff}{aff}
\DeclareMathOperator{\interior}{int}
\DeclareMathOperator{\cl}{cl}
\DeclareMathOperator{\tr}{tr}
\DeclareMathOperator{\diag}{diag}
\DeclareMathOperator{\rank}{rank}
\DeclareMathOperator{\sgn}{sgn}

% --------------------------------------------------------------------------
\begin{document}
\maketitle
\tableofcontents
\newpage

\setcounter{section}{2}

\section{Convex Functions}

\subsection{Definition and Basic Properties}

\begin{definition}[Convex Function]
\label{def:convex-function}
A function $f : \R^n \to \R \cup \{+\infty\}$ is \emph{convex} if $\dom f$ is a convex
set and for all $x, y \in \dom f$ and $\theta \in [0,1]$:
\[
  f(\theta x + (1-\theta) y) \le \theta f(x) + (1-\theta) f(y).
\]
The function is \emph{strictly convex} if the inequality is strict whenever $x \ne y$
and $\theta \in (0,1)$. It is \emph{concave} if $-f$ is convex.
\end{definition}

\begin{definition}[Epigraph]
\label{def:epigraph}
The \emph{epigraph} of $f : \R^n \to \R \cup \{+\infty\}$ is
\[
  \epi f = \{(x, t) \in \R^{n+1} : x \in \dom f, \; f(x) \le t\}.
\]
\end{definition}

\begin{proposition}[Epigraph Characterization]
\label{prop:epigraph}
A function $f$ is convex if and only if $\epi f$ is a convex set.
\end{proposition}

\begin{proof}
($\Rightarrow$) Let $(x_1, t_1), (x_2, t_2) \in \epi f$ and $\theta \in [0,1]$. Then
$x_1, x_2 \in \dom f$, $f(x_1) \le t_1$, $f(x_2) \le t_2$. By convexity of $\dom f$,
$\theta x_1 + (1-\theta) x_2 \in \dom f$, and
\[
  f(\theta x_1 + (1-\theta) x_2) \le \theta f(x_1) + (1-\theta) f(x_2)
  \le \theta t_1 + (1-\theta) t_2.
\]
Hence $(\theta x_1 + (1-\theta) x_2, \; \theta t_1 + (1-\theta) t_2) \in \epi f$.

($\Leftarrow$) Suppose $\epi f$ is convex. For $x_1, x_2 \in \dom f$ and
$\theta \in [0,1]$: $(x_1, f(x_1))$ and $(x_2, f(x_2))$ lie in $\epi f$, so their
convex combination $(\theta x_1 + (1-\theta) x_2, \; \theta f(x_1) + (1-\theta) f(x_2))$
lies in $\epi f$. This means $\theta x_1 + (1-\theta) x_2 \in \dom f$ and $f(\theta x_1
+ (1-\theta) x_2) \le \theta f(x_1) + (1-\theta) f(x_2)$.
\end{proof}

\begin{example}[Standard Convex Functions]
\label{ex:standard-convex}
\begin{enumerate}[label=(\alph*)]
  \item \textbf{Affine}: $f(x) = a^T x + b$ is both convex and concave.
  \item \textbf{Norms}: Every norm $\|\cdot\|$ on $\R^n$ is convex (by the triangle
        inequality and homogeneity).
  \item \textbf{Quadratic}: $f(x) = x^T P x + q^T x + r$ is convex if and only if
        $P \succeq 0$.
  \item \textbf{Exponential}: $f(x) = e^{ax}$ is convex on $\R$ for any $a$.
  \item \textbf{Negative entropy}: $f(x) = x \log x$ is convex on $\R_{++}$.
  \item \textbf{Log-sum-exp}: $f(x) = \log(\sum_{i=1}^n e^{x_i})$ is convex on $\R^n$.
\end{enumerate}
\end{example}

\subsection{Restriction to a Line}

\begin{proposition}[Line Restriction]
\label{prop:line-restriction}
A function $f : \R^n \to \R \cup \{+\infty\}$ is convex if and only if for every
$x \in \dom f$ and every $v \in \R^n$, the function $g(t) = f(x + tv)$ is convex (on
its domain $\{t : x + tv \in \dom f\}$).
\end{proposition}

\begin{proof}
($\Rightarrow$) For $t_1, t_2$ in the domain of $g$ and $\theta \in [0,1]$:
\[
  g(\theta t_1 + (1-\theta) t_2) = f(\theta(x + t_1 v) + (1-\theta)(x + t_2 v))
  \le \theta f(x + t_1 v) + (1-\theta) f(x + t_2 v) = \theta g(t_1) + (1-\theta) g(t_2).
\]
($\Leftarrow$) For $x_1, x_2 \in \dom f$, set $x = x_1$, $v = x_2 - x_1$, so
$g(t) = f(x_1 + t(x_2 - x_1))$, with $g(0) = f(x_1)$, $g(1) = f(x_2)$. Convexity
of $g$ gives $g(\theta) \le \theta g(1) + (1-\theta) g(0) = \theta f(x_2) +
(1-\theta) f(x_1)$.
\end{proof}

\subsection{First-Order Condition}

\begin{theorem}[First-Order Condition for Convexity]
\label{thm:first-order}
Let $f : \R^n \to \R$ be differentiable with convex domain. Then $f$ is convex if and
only if for all $x, y \in \dom f$:
\begin{equation}
  \label{eq:first-order}
  f(y) \ge f(x) + \nabla f(x)^T (y - x).
\end{equation}
\end{theorem}

\begin{proof}
($\Rightarrow$) Assume $f$ is convex. For $x, y \in \dom f$ and $t \in (0, 1]$:
\[
  f(x + t(y - x)) \le (1-t) f(x) + t f(y).
\]
Rearranging: $f(y) \ge f(x) + \frac{f(x + t(y-x)) - f(x)}{t}$. Taking $t \to 0^+$, the
right-hand side becomes $f(x) + \nabla f(x)^T(y - x)$ by the definition of the
directional derivative.

($\Leftarrow$) Assume \eqref{eq:first-order} holds for all $x, y$. Let $x_1, x_2 \in
\dom f$, $\theta \in [0,1]$, and $z = \theta x_1 + (1-\theta) x_2$. Apply
\eqref{eq:first-order} with $x = z$:
\begin{align*}
  f(x_1) &\ge f(z) + \nabla f(z)^T(x_1 - z), \\
  f(x_2) &\ge f(z) + \nabla f(z)^T(x_2 - z).
\end{align*}
Multiply the first by $\theta$ and the second by $(1-\theta)$ and add:
\[
  \theta f(x_1) + (1-\theta) f(x_2) \ge f(z) +
  \nabla f(z)^T \bigl(\theta(x_1 - z) + (1-\theta)(x_2 - z)\bigr) = f(z),
\]
since $\theta x_1 + (1-\theta) x_2 - z = 0$.
\end{proof}

\begin{keyresult}
The first-order condition \eqref{eq:first-order} says the tangent hyperplane at any point
is a \emph{global underestimator} of a convex function. This has a powerful consequence:
if $\nabla f(x^*) = 0$, then $f(y) \ge f(x^*)$ for all $y$, so $x^*$ is a global
minimizer.
\end{keyresult}

\begin{rlconnection}
The first-order condition directly implies that \textbf{any local minimum of a convex
function is a global minimum}. In RL, this is why convex surrogate losses (e.g., the
linear approximation in TRPO) can be minimized greedily without fear of local optima.
For non-convex objectives (deep RL), this guarantee fails, motivating trust region and
proximal methods to control the update step.
\end{rlconnection}

\subsection{Second-Order Condition}

\begin{theorem}[Second-Order Condition for Convexity]
\label{thm:second-order}
Let $f : \R^n \to \R$ be twice differentiable with open convex domain. Then $f$ is convex
if and only if the Hessian is positive semidefinite everywhere:
\[
  \nabla^2 f(x) \succeq 0 \quad \text{for all } x \in \dom f.
\]
$f$ is strictly convex if $\nabla^2 f(x) \succ 0$ for all $x$ (the converse does not
hold: $f(x) = x^4$ is strictly convex but $f''(0) = 0$).
\end{theorem}

\begin{proof}
By the line restriction (Proposition~\ref{prop:line-restriction}), it suffices to prove
the one-dimensional case: $g : \R \to \R$ twice differentiable is convex if and only if
$g''(t) \ge 0$ for all $t$.

($\Rightarrow$) Suppose $g$ is convex and $g''(t_0) < 0$ for some $t_0$. By the
first-order condition:
\[
  g(t_0 + h) \ge g(t_0) + g'(t_0) h \quad \text{for all } h.
\]
By Taylor's theorem, $g(t_0 + h) = g(t_0) + g'(t_0) h + \tfrac{1}{2} g''(\xi) h^2$ for
some $\xi$ between $t_0$ and $t_0 + h$. For sufficiently small $|h|$, $g''(\xi) < 0$ by
continuity, so $g(t_0 + h) < g(t_0) + g'(t_0) h$, contradicting the first-order
condition.

($\Leftarrow$) Suppose $g''(t) \ge 0$ for all $t$. By Taylor's theorem, for any $x, y$:
\[
  g(y) = g(x) + g'(x)(y - x) + \tfrac{1}{2} g''(\xi)(y - x)^2
  \ge g(x) + g'(x)(y - x),
\]
which is the first-order condition. By Theorem~\ref{thm:first-order}, $g$ is convex.

For the multivariate case: $f$ is convex iff $g(t) = f(x + tv)$ is convex for all $x, v$.
We have $g''(t) = v^T \nabla^2 f(x + tv) v$. Thus $g''(t) \ge 0$ for all $t, v$ iff
$\nabla^2 f \succeq 0$ everywhere.
\end{proof}

\begin{example}[Quadratic Functions]
\label{ex:quadratic-convexity}
For $f(x) = \tfrac{1}{2} x^T P x + q^T x + r$, we have $\nabla^2 f(x) = P$. So $f$ is
convex iff $P \succeq 0$, and strictly convex iff $P \succ 0$.
\end{example}

\begin{example}[Log-Sum-Exp]
\label{ex:lse-convex}
For $f(x) = \log \sum_{i=1}^n e^{x_i}$, the Hessian is
\[
  \nabla^2 f(x) = \frac{1}{\mathbf{1}^T z} \diag(z)
    - \frac{1}{(\mathbf{1}^T z)^2} z z^T,
\]
where $z_i = e^{x_i}$. For any $v$:
\[
  v^T \nabla^2 f(x) v = \frac{(\sum z_i)(\sum z_i v_i^2) - (\sum z_i v_i)^2}
  {(\sum z_i)^2} \ge 0,
\]
by the Cauchy--Schwarz inequality (with weights $z_i$). Hence $f$ is convex.
\end{example}

\subsection{Jensen's Inequality}

\begin{theorem}[Jensen's Inequality]
\label{thm:jensen}
Let $f : \R^n \to \R$ be convex and let $X$ be a random vector with $\E[\|X\|] < \infty$
and $\E[X] \in \dom f$. Then
\[
  f(\E[X]) \le \E[f(X)].
\]
If $f$ is concave, the inequality reverses.
\end{theorem}

\begin{proof}
By the first-order condition (Theorem~\ref{thm:first-order}) applied at $x = \E[X]$:
\[
  f(X) \ge f(\E[X]) + \nabla f(\E[X])^T (X - \E[X]).
\]
Taking expectations:
\[
  \E[f(X)] \ge f(\E[X]) + \nabla f(\E[X])^T (\E[X] - \E[X]) = f(\E[X]).
\]
For the non-differentiable case, use the existence of a subgradient $g \in \partial
f(\E[X])$ satisfying $f(y) \ge f(\E[X]) + g^T(y - \E[X])$ for all $y$ (subgradients
exist for all convex functions at interior points of their domain). The same argument
applies.
\end{proof}

\begin{keyresult}[title={Jensen's Inequality -- Finite Form}]
For convex $f$, points $x_1, \ldots, x_k$, and weights $\theta_i \ge 0$ with
$\sum \theta_i = 1$:
\[
  f\!\left(\sum_{i=1}^k \theta_i x_i\right) \le \sum_{i=1}^k \theta_i f(x_i).
\]
This is the defining property of convexity extended from two points to $k$ points by
induction.
\end{keyresult}

\begin{rlconnection}
Jensen's inequality is ubiquitous in RL/ML analysis:
\begin{itemize}
  \item \textbf{Evidence lower bound (ELBO)}: Since $\log$ is concave,
        $\log \E[f(X)] \ge \E[\log f(X)]$, which is the foundation of variational
        inference.
  \item \textbf{KL divergence}: Non-negativity of $D_{\text{KL}}(p\|q) =
        \E_p[\log(p/q)] \ge 0$ follows from Jensen applied to $-\log$ (convex).
  \item \textbf{Importance sampling}: Variance analysis of importance-weighted
        estimators uses Jensen to bound $\E[(w(X) f(X))^2]$ terms.
\end{itemize}
\end{rlconnection}

%% ========================================================================
%% SECTION 4: OPERATIONS PRESERVING CONVEXITY OF FUNCTIONS
%% ========================================================================

\section{Operations Preserving Convexity of Functions}

\subsection{Nonnegative Weighted Sums}

\begin{proposition}[Nonnegative Weighted Sum]
\label{prop:nonneg-sum}
If $f_1, \ldots, f_m$ are convex and $w_1, \ldots, w_m \ge 0$, then $f = \sum_{i=1}^m
w_i f_i$ is convex.
\end{proposition}

\begin{proof}
For $\theta \in [0,1]$:
\[
  f(\theta x + (1-\theta) y) = \sum_i w_i f_i(\theta x + (1-\theta)y)
  \le \sum_i w_i [\theta f_i(x) + (1-\theta) f_i(y)]
  = \theta f(x) + (1-\theta) f(y). \qedhere
\]
\end{proof}

\subsection{Composition with Affine Maps}

\begin{proposition}[Affine Composition]
\label{prop:affine-comp}
If $f : \R^m \to \R$ is convex, $A \in \R^{m \times n}$, and $b \in \R^m$, then
$g(x) = f(Ax + b)$ is convex.
\end{proposition}

\begin{proof}
$\dom g = \{x : Ax + b \in \dom f\}$ is convex (preimage of convex set under affine
map). For $\theta \in [0,1]$:
\begin{align*}
  g(\theta x + (1-\theta) y) &= f(A(\theta x + (1-\theta)y) + b) \\
  &= f(\theta(Ax+b) + (1-\theta)(Ay+b)) \\
  &\le \theta f(Ax+b) + (1-\theta) f(Ay+b) = \theta g(x) + (1-\theta) g(y). \qedhere
\end{align*}
\end{proof}

\subsection{Pointwise Maximum and Supremum}

\begin{proposition}[Pointwise Maximum]
\label{prop:pointwise-max}
If $f_1, \ldots, f_m$ are convex, then $f(x) = \max\{f_1(x), \ldots, f_m(x)\}$ is
convex.
\end{proposition}

\begin{proof}
For $\theta \in [0,1]$:
\[
  f(\theta x + (1-\theta)y) = \max_i f_i(\theta x + (1-\theta) y)
  \le \max_i [\theta f_i(x) + (1-\theta) f_i(y)]
  \le \theta \max_i f_i(x) + (1-\theta) \max_i f_i(y).
\]
The last inequality uses: $\max_i (a_i + b_i) \le \max_i a_i + \max_i b_i$ (applied
with $a_i = \theta f_i(x)$ and $b_i = (1-\theta) f_i(y)$, noting $\theta \ge 0$ and
$1-\theta \ge 0$).
\end{proof}

\begin{proposition}[Pointwise Supremum]
\label{prop:pointwise-sup}
If $f(x, y)$ is convex in $x$ for each $y \in \mathcal{A}$, then
$g(x) = \sup_{y \in \mathcal{A}} f(x, y)$ is convex (with $\dom g = \{x :
g(x) < +\infty\}$).
\end{proposition}

\begin{proof}
The epigraph of $g$ is $\epi g = \bigcap_{y \in \mathcal{A}} \epi f(\cdot, y)$. Each
$\epi f(\cdot, y)$ is convex (since $f(\cdot, y)$ is convex), and the intersection of
convex sets is convex (Proposition~\ref{prop:intersection}). Hence $\epi g$ is convex,
so $g$ is convex (Proposition~\ref{prop:epigraph}).
\end{proof}

\begin{rlconnection}
The pointwise supremum rule shows why the \textbf{support function}
$S_C(y) = \sup_{x \in C} y^T x$ is convex for any set $C$ (not necessarily convex).
In RL, the optimal value function $V^*(s) = \sup_\pi V^\pi(s)$ is the supremum of
affine functions of the transition model parameters in certain formulations, connecting
to the convexity of the value function in the occupancy measure.
\end{rlconnection}

\subsection{Composition Rules}

\begin{proposition}[Scalar Composition]
\label{prop:scalar-composition}
Let $h : \R \to \R$ and $g : \R^n \to \R$. Define $f = h \circ g$. Then:
\begin{enumerate}[label=(\alph*)]
  \item $f$ is convex if $h$ is convex and nondecreasing, and $g$ is convex.
  \item $f$ is convex if $h$ is convex and nonincreasing, and $g$ is concave.
  \item $f$ is concave if $h$ is concave and nondecreasing, and $g$ is concave.
  \item $f$ is concave if $h$ is concave and nonincreasing, and $g$ is convex.
\end{enumerate}
\end{proposition}

\begin{proof}
We prove (a); the others are analogous. Assume $h, g$ twice differentiable for
clarity. Then $f''(x) = h''(g(x))(g'(x))^2 + h'(g(x)) g''(x)$ (in 1D). Since $h$ is
convex, $h'' \ge 0$; since $h$ is nondecreasing, $h' \ge 0$; since $g$ is convex,
$g'' \ge 0$. Both terms are nonnegative, so $f'' \ge 0$.

For the general (non-differentiable) proof of (a): let $\theta \in [0,1]$. Since $g$
is convex: $g(\theta x + (1-\theta) y) \le \theta g(x) + (1-\theta) g(y)$. Since $h$
is nondecreasing:
\[
  h(g(\theta x + (1-\theta) y)) \le h(\theta g(x) + (1-\theta) g(y)).
\]
Since $h$ is convex:
\[
  h(\theta g(x) + (1-\theta) g(y)) \le \theta h(g(x)) + (1-\theta) h(g(y)).
\]
Combining: $f(\theta x + (1-\theta)y) \le \theta f(x) + (1-\theta) f(y)$.
\end{proof}

\begin{example}[Applications of Composition Rules]
\label{ex:composition-apps}
\begin{enumerate}[label=(\alph*)]
  \item $e^{g(x)}$ is convex if $g$ is convex ($h(t) = e^t$ is convex and
        nondecreasing).
  \item $1/g(x)$ is convex if $g$ is concave and positive ($h(t) = 1/t$ is convex and
        nonincreasing on $\R_{++}$).
  \item $\log g(x)$ is concave if $g$ is concave and positive ($h(t) = \log t$ is
        concave and nondecreasing).
  \item $\|Ax + b\|_p$ is convex ($h = \|\cdot\|_p$ is convex and nondecreasing on
        the nonnegative orthant is not needed here; use affine composition
        Proposition~\ref{prop:affine-comp} directly since norms are convex).
\end{enumerate}
\end{example}

\subsection{Infimum / Minimization}

\begin{proposition}[Partial Minimization]
\label{prop:partial-min}
If $f(x, y)$ is convex in $(x, y)$ and $C$ is a convex set, then
$g(x) = \inf_{y \in C} f(x, y)$ is convex (provided $g(x) > -\infty$ for all $x$).
\end{proposition}

\begin{proof}
Let $x_1, x_2 \in \dom g$ and $\theta \in [0,1]$. For any $\varepsilon > 0$, choose
$y_1, y_2 \in C$ such that $f(x_i, y_i) \le g(x_i) + \varepsilon$. Then
\begin{align*}
  g(\theta x_1 + (1-\theta) x_2)
  &\le f(\theta x_1 + (1-\theta)x_2, \; \theta y_1 + (1-\theta) y_2) \\
  &\le \theta f(x_1, y_1) + (1-\theta) f(x_2, y_2) \\
  &\le \theta g(x_1) + (1-\theta) g(x_2) + \varepsilon.
\end{align*}
Since $\varepsilon > 0$ is arbitrary, $g(\theta x_1 + (1-\theta) x_2) \le \theta g(x_1)
+ (1-\theta) g(x_2)$.
\end{proof}

\subsection{Perspective Transform}

\begin{proposition}[Perspective of a Convex Function]
\label{prop:perspective-function}
If $f : \R^n \to \R$ is convex, then its \emph{perspective}
$g : \R^{n+1} \to \R$ defined by $g(x, t) = t f(x/t)$ with $\dom g = \{(x,t) :
x/t \in \dom f, \; t > 0\}$ is convex.
\end{proposition}

\begin{proof}
We show $\epi g$ is convex. Note $(x, t, s) \in \epi g$ iff $t > 0$ and
$tf(x/t) \le s$, iff $f(x/t) \le s/t$, iff $(x/t, s/t) \in \epi f$. Thus
$(x, t, s) \in \epi g$ iff $(x/t, s/t) \in \epi f$ and $t > 0$. This is the preimage
of $\epi f$ under the perspective map $(x, t, s) \mapsto (x/t, s/t)$.
By Proposition~\ref{prop:perspective}, the perspective map preserves convexity of
preimages. Since $\epi f$ is convex (as $f$ is convex), $\epi g$ is convex.
\end{proof}

\begin{example}[KL Divergence from Perspective]
\label{ex:kl-perspective}
The negative entropy $f(x) = x \log x$ is convex on $\R_{++}$. Its perspective is
$g(x, t) = t \cdot (x/t) \log(x/t) = x \log(x/t)$. Summing over coordinates gives the
KL divergence: $D_{\text{KL}}(p \| q) = \sum_i p_i \log(p_i / q_i)$, which is jointly
convex in $(p, q)$.
\end{example}

%% ========================================================================
%% SECTION 5: CONJUGATE FUNCTIONS
%% ========================================================================

\section{Conjugate Functions}

\subsection{The Fenchel Conjugate}

\begin{definition}[Fenchel Conjugate]
\label{def:conjugate}
Let $f : \R^n \to \R \cup \{+\infty\}$. The \emph{Fenchel conjugate} (or
\emph{convex conjugate}) of $f$ is $f^* : \R^n \to \R \cup \{+\infty\}$ defined by
\[
  f^*(y) = \sup_{x \in \dom f} \left( y^T x - f(x) \right).
\]
\end{definition}

\begin{proposition}[Conjugate is Always Convex]
\label{prop:conjugate-convex}
For any function $f$ (not necessarily convex), $f^*$ is convex and closed
(i.e., lower semicontinuous).
\end{proposition}

\begin{proof}
For each fixed $x$, the function $y \mapsto y^T x - f(x)$ is affine in $y$, hence
convex. The conjugate $f^*(y) = \sup_x (y^T x - f(x))$ is the pointwise supremum of
convex (in fact affine) functions of $y$. By Proposition~\ref{prop:pointwise-sup},
$f^*$ is convex. Closedness follows because $\epi f^* = \bigcap_x \{(y,t) :
y^T x - f(x) \le t\}$ is an intersection of closed halfspaces.
\end{proof}

\begin{theorem}[Fenchel's Inequality]
\label{thm:fenchel-inequality}
For any function $f$ and any $x, y \in \R^n$:
\[
  f(x) + f^*(y) \ge y^T x.
\]
\end{theorem}

\begin{proof}
By definition, $f^*(y) = \sup_z (y^T z - f(z)) \ge y^T x - f(x)$. Rearranging gives
$f(x) + f^*(y) \ge y^T x$.
\end{proof}

\begin{keyresult}[title={Fenchel's Inequality}]
The inequality $f(x) + f^*(y) \ge \langle y, x \rangle$ is often called the
\emph{Young--Fenchel inequality}. Equality holds if and only if $y \in \partial f(x)$
(the subdifferential of $f$ at $x$), when $f$ is convex and closed.
\end{keyresult}

\subsection{Conjugate of the Conjugate}

\begin{theorem}[Biconjugate Theorem]
\label{thm:biconjugate}
If $f : \R^n \to \R \cup \{+\infty\}$ is convex and closed (lower semicontinuous) with
$f \not\equiv +\infty$, then $f^{**} = f$.
\end{theorem}

\begin{proof}[Proof sketch]
By Fenchel's inequality, $f^{**}(x) = \sup_y (x^T y - f^*(y)) \le f(x)$ for all $x$
(set $z = x$ in the supremum defining $f^*$, then rearrange). The reverse inequality
$f^{**} \ge f$ uses the supporting hyperplane theorem: since $\epi f$ is a closed convex
set and $(x, f(x))$ lies on its boundary, there is a supporting hyperplane
$a^T u + bt \ge a^T x + b f(x)$ for all $(u, t) \in \epi f$. Analyzing this gives
$f^*(-a/b) = (-a/b)^T x - f(x)$, hence $f^{**}(x) \ge (-a/b)^T x - f^*(-a/b) = f(x)$.
A complete proof requires careful handling of the case $b = 0$ and uses the full
Fenchel--Moreau theorem from convex analysis.
\end{proof}

\subsection{Examples of Conjugates}

\begin{example}[Conjugate of a Linear Function]
\label{ex:conjugate-linear}
Let $f(x) = a^T x + b$. Then
\[
  f^*(y) = \sup_x (y^T x - a^T x - b) = \sup_x ((y - a)^T x) - b
  = \begin{cases} -b & \text{if } y = a, \\ +\infty & \text{otherwise.} \end{cases}
\]
\end{example}

\begin{example}[Conjugate of a Quadratic]
\label{ex:conjugate-quadratic}
Let $f(x) = \tfrac{1}{2} x^T Q x$ where $Q \succ 0$. Then
\[
  f^*(y) = \sup_x \left( y^T x - \tfrac{1}{2} x^T Q x \right).
\]
Setting the gradient to zero: $y - Qx = 0$, so $x^* = Q^{-1} y$. Substituting:
\[
  f^*(y) = y^T Q^{-1} y - \tfrac{1}{2} y^T Q^{-1} Q Q^{-1} y
  = \tfrac{1}{2} y^T Q^{-1} y.
\]
\end{example}

\begin{example}[Conjugate of the Negative Entropy]
\label{ex:conjugate-entropy}
Let $f(x) = x \log x$ for $x > 0$. Then
$f^*(y) = \sup_{x > 0} (yx - x \log x)$. Setting the derivative to zero: $y - \log x
- 1 = 0$, so $x^* = e^{y-1}$. Thus $f^*(y) = e^{y-1}$.
\end{example}

\begin{example}[Conjugate of Indicator Function]
\label{ex:conjugate-indicator}
Let $\delta_C(x) = 0$ if $x \in C$, $+\infty$ otherwise. Then
\[
  \delta_C^*(y) = \sup_{x \in C} y^T x = S_C(y),
\]
the \emph{support function} of $C$. This connects set geometry to duality.
\end{example}

\begin{example}[Conjugate of Norm]
\label{ex:conjugate-norm}
Let $f(x) = \|x\|$. Then $f^*(y) = \delta_{\{y : \|y\|_* \le 1\}}(y)$, where
$\|y\|_* = \sup_{\|x\| \le 1} y^T x$ is the dual norm. To see this: $f^*(y) =
\sup_x (y^T x - \|x\|)$. If $\|y\|_* > 1$, then there exists $x_0$ with $\|x_0\| = 1$
and $y^T x_0 > 1$, so $f^*(y) \ge \sup_{t \ge 0} (t y^T x_0 - t) = +\infty$. If
$\|y\|_* \le 1$, then $y^T x \le \|y\|_* \|x\| \le \|x\|$, so $f^*(y) \le 0$; and
$f^*(y) \ge y^T \cdot 0 - \|0\| = 0$. Thus $f^*(y) = 0$.
\end{example}

\begin{rlconnection}
Fenchel conjugates appear naturally in \textbf{regularized MDPs}. If the policy is
regularized by an entropy term $\Omega(\pi) = \sum_a \pi(a) \log \pi(a)$, the
regularized Bellman operator involves $\Omega^*$, the conjugate of the entropy, which
equals the log-sum-exp function (the ``soft maximum''). This is the foundation of
\textbf{soft Q-learning} and \textbf{SAC (Soft Actor-Critic)}, where the softmax
policy arises as the conjugate dual solution.
\end{rlconnection}

%% ========================================================================
%% SECTION 6: SUBLEVEL SETS AND QUASICONVEX FUNCTIONS
%% ========================================================================

\section{Sublevel Sets and Quasiconvex Functions}

\subsection{Sublevel Sets}

\begin{definition}[Sublevel Set]
\label{def:sublevel}
The $\alpha$-\emph{sublevel set} of $f : \R^n \to \R$ is
$C_\alpha = \{x \in \dom f : f(x) \le \alpha\}$.
\end{definition}

\begin{proposition}[Sublevel Sets of Convex Functions]
\label{prop:sublevel}
If $f$ is convex, then all its sublevel sets are convex.
\end{proposition}

\begin{proof}
Let $x, y \in C_\alpha$ (so $f(x) \le \alpha$ and $f(y) \le \alpha$) and
$\theta \in [0,1]$. Then
\[
  f(\theta x + (1-\theta) y) \le \theta f(x) + (1-\theta) f(y)
  \le \theta \alpha + (1-\theta) \alpha = \alpha,
\]
so $\theta x + (1-\theta) y \in C_\alpha$.
\end{proof}

\begin{warning}
The converse of Proposition~\ref{prop:sublevel} is \textbf{false}. A function can have
all sublevel sets convex without being convex. For example, $f(x) = -e^x$ has convex
sublevel sets $\{x : -e^x \le \alpha\} = \{x : x \ge \log(-\alpha)\}$ (halflines)
for $\alpha < 0$, but $f$ is concave (not convex). Functions with convex sublevel sets
are called \emph{quasiconvex} (see below).
\end{warning}

\subsection{Quasiconvex Functions}

\begin{definition}[Quasiconvex Function]
\label{def:quasiconvex}
A function $f : \R^n \to \R$ is \emph{quasiconvex} if $\dom f$ is convex and all
sublevel sets $\{x : f(x) \le \alpha\}$ are convex. Equivalently, for all
$x, y \in \dom f$ and $\theta \in [0,1]$:
\[
  f(\theta x + (1-\theta) y) \le \max\{f(x), f(y)\}.
\]
$f$ is \emph{quasiconcave} if $-f$ is quasiconvex, i.e., all \emph{superlevel sets}
$\{x : f(x) \ge \alpha\}$ are convex.
\end{definition}

\begin{proposition}[Equivalence of Quasiconvexity Definitions]
\label{prop:quasiconvex-equiv}
$f$ has all sublevel sets convex if and only if $f(\theta x + (1-\theta)y) \le
\max\{f(x), f(y)\}$ for all $x, y \in \dom f$ and $\theta \in [0,1]$.
\end{proposition}

\begin{proof}
($\Rightarrow$) Let $\alpha = \max\{f(x), f(y)\}$. Then $x, y \in C_\alpha$, which is
convex, so $\theta x + (1-\theta) y \in C_\alpha$, giving
$f(\theta x + (1-\theta) y) \le \alpha = \max\{f(x), f(y)\}$.

($\Leftarrow$) Fix $\alpha$ and let $x, y \in C_\alpha$. Then $\max\{f(x), f(y)\} \le
\alpha$, so $f(\theta x + (1-\theta) y) \le \alpha$, giving
$\theta x + (1-\theta) y \in C_\alpha$.
\end{proof}

\begin{example}[Quasiconvex Functions]
\label{ex:quasiconvex}
\begin{enumerate}[label=(\alph*)]
  \item Every convex function is quasiconvex.
  \item $f(x) = \lceil x \rceil$ (ceiling function) is quasiconvex.
  \item $\log x$ is both quasiconvex and quasiconcave on $\R_{++}$ (it is monotone, and
        monotone functions are both quasiconvex and quasiconcave).
  \item The ratio $f(x) = p(x)/q(x)$ where $p$ is convex nonneg and $q$ is concave
        positive is quasiconvex.
\end{enumerate}
\end{example}

\begin{proposition}[First-Order Condition for Quasiconvexity]
\label{prop:quasiconvex-first-order}
A differentiable function $f$ with convex domain is quasiconvex if and only if for all
$x, y \in \dom f$:
\[
  f(y) \le f(x) \implies \nabla f(x)^T(y - x) \le 0.
\]
\end{proposition}

\begin{proof}
($\Rightarrow$) Suppose $f$ is quasiconvex and $f(y) \le f(x)$. For $t \in (0,1]$:
$f(x + t(y-x)) \le \max\{f(x), f(y)\} = f(x)$. So $\frac{f(x+t(y-x)) - f(x)}{t}
\le 0$. Taking $t \to 0^+$: $\nabla f(x)^T(y - x) \le 0$.

($\Leftarrow$) We prove the contrapositive. Suppose $f$ is not quasiconvex, so there
exist $x, y$ and $\theta \in (0,1)$ with $f(z) > \max\{f(x), f(y)\}$ where
$z = \theta x + (1-\theta) y$. Then $f(x) < f(z)$ and $f(y) < f(z)$. The condition
would require $\nabla f(z)^T(x - z) \le 0$ and $\nabla f(z)^T(y - z) \le 0$. But
$\theta(x-z) + (1-\theta)(y-z) = 0$, and both inner products nonpositive with a
nontrivial linear combination equal to zero forces both to be zero. A more careful
argument using the mean value theorem on the line segment shows a contradiction with
the strict inequality $f(z) > f(x)$.
\end{proof}

%% ========================================================================
%% SECTION 7: LOG-CONCAVE AND LOG-CONVEX FUNCTIONS
%% ========================================================================

\section{Log-Concave and Log-Convex Functions}

\begin{definition}[Log-Concave and Log-Convex Functions]
\label{def:log-concave}
A function $f : \R^n \to \R$ with $f(x) > 0$ for all $x \in \dom f$ is
\emph{log-concave} if $\log f$ is concave, i.e., for all $x, y \in \dom f$ and
$\theta \in [0,1]$:
\[
  f(\theta x + (1-\theta) y) \ge f(x)^\theta f(y)^{1-\theta}.
\]
$f$ is \emph{log-convex} if $\log f$ is convex, equivalently:
\[
  f(\theta x + (1-\theta) y) \le f(x)^\theta f(y)^{1-\theta}.
\]
\end{definition}

\begin{remark}
The defining inequality $f(\theta x + (1-\theta) y) \ge f(x)^\theta f(y)^{1-\theta}$
follows from exponentiating the concavity of $\log f$:
\[
  \log f(\theta x + (1-\theta) y) \ge \theta \log f(x) + (1-\theta) \log f(y)
  = \log\bigl(f(x)^\theta f(y)^{1-\theta}\bigr).
\]
\end{remark}

\begin{proposition}[Hierarchy of Convexity Notions]
\label{prop:hierarchy}
For a positive function $f$:
\begin{center}
  log-convex $\implies$ convex $\implies$ quasiconvex, \\[4pt]
  concave $\implies$ log-concave $\implies$ quasiconcave.
\end{center}
\end{proposition}

\begin{proof}
\textbf{Log-convex $\implies$ convex}: If $\log f$ is convex, then $f = e^{\log f}$ is
the composition of the convex nondecreasing function $e^t$ with the convex function
$\log f$, hence convex by Proposition~\ref{prop:scalar-composition}(a).

\textbf{Convex $\implies$ quasiconvex}: Already proved (sublevel sets of convex
functions are convex).

\textbf{Concave $\implies$ log-concave}: We must show $\log f$ is concave, assuming $f$ is
concave and positive. By the AM--GM inequality, for $\theta \in [0,1]$:
\[
  f(\theta x + (1-\theta) y) \ge \theta f(x) + (1-\theta) f(y)
  \ge f(x)^\theta f(y)^{1-\theta},
\]
where the first inequality is concavity of $f$ and the second is the weighted AM--GM
inequality. Taking logarithms gives concavity of $\log f$.

\textbf{Log-concave $\implies$ quasiconcave}: If $\log f$ is concave, then $\log f$
has convex superlevel sets (since concave functions have convex superlevel sets).
The $\alpha$-superlevel set of $f$ is $\{x : f(x) \ge \alpha\} = \{x : \log f(x) \ge
\log \alpha\}$ (for $\alpha > 0$), which is a superlevel set of the concave function
$\log f$, hence convex.
\end{proof}

\begin{example}[Log-Concave Distributions]
\label{ex:log-concave-dist}
Many common probability densities are log-concave:
\begin{enumerate}[label=(\alph*)]
  \item \textbf{Gaussian}: $f(x) \propto \exp(-\tfrac{1}{2}(x-\mu)^T \Sigma^{-1}
        (x-\mu))$. Here $\log f$ is a concave quadratic.
  \item \textbf{Exponential}: $f(x) = \lambda e^{-\lambda x}$ for $x \ge 0$.
        $\log f = \log \lambda - \lambda x$ is concave (affine).
  \item \textbf{Uniform}: $f(x) = 1/(b-a)$ on $[a,b]$. $\log f$ is constant (concave).
  \item \textbf{Laplace}: $f(x) \propto e^{-|x-\mu|/b}$. $\log f$ is concave (piecewise
        linear with negative slopes).
\end{enumerate}
\end{example}

\begin{proposition}[Properties of Log-Concave Functions]
\label{prop:log-concave-properties}
\begin{enumerate}[label=(\alph*)]
  \item Products: if $f, g$ are log-concave, then $fg$ is log-concave.
  \item Marginals: if $f(x, y)$ is log-concave, then $g(x) = \int f(x, y) \, dy$ is
        log-concave (Prekopa's theorem).
  \item Convolution: if $f, g$ are log-concave, then $(f * g)(x) = \int f(x-y) g(y) \,
        dy$ is log-concave.
\end{enumerate}
\end{proposition}

\begin{proof}
(a) $\log(fg) = \log f + \log g$ is the sum of two concave functions, hence concave.

(b) This is Prekopa's theorem. The proof uses the Prekopa--Leindler inequality: if $f, g,
h : \R^n \to \R_{\ge 0}$ satisfy $h(\theta x + (1-\theta)y) \ge f(x)^\theta
g(y)^{1-\theta}$ for all $x, y$ and some $\theta \in (0,1)$, then
\[
  \int h \ge \left(\int f\right)^\theta \left(\int g\right)^{1-\theta}.
\]
Apply this to the ``slices'' of the log-concave function in the variable being
integrated out. Specifically, for fixed $x_1, x_2$ and $\bar{x} = \theta x_1 +
(1-\theta) x_2$, the log-concavity of $f$ gives
$f(\bar{x}, y) \ge f(x_1, y_1)^\theta f(x_2, y_2)^{1-\theta}$ when
$y = \theta y_1 + (1-\theta) y_2$. Applying the Prekopa--Leindler inequality in $y$:
$g(\bar{x}) \ge g(x_1)^\theta g(x_2)^{1-\theta}$.

(c) Write $(f * g)(x) = \int f(x - y) g(y) \, dy = \int h(x, y) \, dy$ where
$h(x, y) = f(x - y) g(y)$. The function $h$ is log-concave in $(x, y)$ since
$\log h(x,y) = \log f(x-y) + \log g(y)$ is the sum of a concave function of $(x,y)$
(composition of concave $\log f$ with affine $(x,y) \mapsto x - y$) and a concave
function of $y$ (hence of $(x,y)$). By part (b), the integral over $y$ is log-concave
in $x$.
\end{proof}

\begin{rlconnection}
Log-concavity is central to Bayesian inference in machine learning:
\begin{itemize}
  \item A Gaussian prior times a log-concave likelihood gives a log-concave posterior
        (by the product rule), guaranteeing unimodality and enabling efficient sampling
        (e.g., Langevin dynamics converges in polynomial time for log-concave targets).
  \item The marginal property ensures that marginalizing latent variables from a
        log-concave joint preserves log-concavity, important for evidence computation.
  \item In RL, the softmax policy $\pi(a|s) \propto \exp(Q(s,a)/\tau)$ is log-concave
        (in fact log-linear) in $Q$, which simplifies the analysis of policy gradient
        methods.
\end{itemize}
\end{rlconnection}

%% ========================================================================
%% SECTION 8: EXERCISES
%% ========================================================================

\section{Exercises}

\begin{exercise}[Convexity of Set Operations]
\label{ex:set-operations}
Let $C_1, C_2 \subseteq \R^n$ be convex sets. Prove that the following are convex:
\begin{enumerate}[label=(\alph*)]
  \item The Minkowski sum $C_1 + C_2 = \{x_1 + x_2 : x_1 \in C_1, x_2 \in C_2\}$.
  \item The set $\{(\lambda, x) : x \in \lambda C_1, \; \lambda \ge 0\}
        \subseteq \R^{n+1}$.
\end{enumerate}
\end{exercise}

\begin{exercise}[Convexity via Epigraph]
\label{ex:epigraph-exercise}
Use the epigraph characterization to prove that $f(x) = \max\{|x_1|, \ldots, |x_n|\}$
(the $\ell_\infty$-norm) is convex.
\textit{Hint}: Express $\epi f$ as an intersection of halfspaces.
\end{exercise}

\begin{exercise}[First-Order Condition Application]
\label{ex:first-order-app}
Let $f : \R^n \to \R$ be convex and differentiable.
\begin{enumerate}[label=(\alph*)]
  \item Show that $\nabla f(x) = 0$ implies $x$ is a global minimizer.
  \item Show that if $f$ is strictly convex, there is at most one global minimizer.
  \item Give an example of a convex function with no minimizer, a unique minimizer,
        and a non-unique set of minimizers.
\end{enumerate}
\end{exercise}

\begin{exercise}[Conjugate Computation]
\label{ex:conjugate-compute}
Compute the Fenchel conjugate of each function:
\begin{enumerate}[label=(\alph*)]
  \item $f(x) = \|x\|_1$ (the $\ell_1$-norm on $\R^n$).
  \item $f(x) = -\log x$ for $x > 0$.
  \item $f(x) = \frac{1}{p}\|x\|_p^p$ for $p > 1$.
\end{enumerate}
\end{exercise}

\begin{exercise}[Jensen's Inequality Applications]
\label{ex:jensen-apps}
Using Jensen's inequality:
\begin{enumerate}[label=(\alph*)]
  \item Prove $\E[X^2] \ge (\E[X])^2$ for any random variable $X$ (i.e., $\text{Var}(X) \ge 0$).
  \item Prove that for positive weights $w_i$ summing to 1 and positive reals $a_i$:
        $\prod_{i=1}^n a_i^{w_i} \le \sum_{i=1}^n w_i a_i$ (weighted AM--GM).
  \item Prove that $D_{\text{KL}}(p \| q) \ge 0$ for any two probability distributions
        $p, q$ on a finite set.
\end{enumerate}
\end{exercise}

\begin{exercise}[Log-Sum-Exp Properties]
\label{ex:lse-properties}
Let $\text{LSE}(x) = \log(\sum_{i=1}^n e^{x_i})$.
\begin{enumerate}[label=(\alph*)]
  \item Verify that $\max_i x_i \le \text{LSE}(x) \le \max_i x_i + \log n$.
  \item Show that $\text{LSE}$ is the conjugate of the negative entropy restricted to the
        probability simplex: $\text{LSE}(y) = \sup_{x \in \Delta_n}
        (y^T x + \sum_i x_i \log x_i)$ where $\Delta_n$ is the simplex.
        \textit{Hint}: Use Lagrange multipliers.
\end{enumerate}
\end{exercise}

\begin{exercise}[Perspective Transform and KL Divergence]
\label{ex:perspective-kl}
\begin{enumerate}[label=(\alph*)]
  \item Verify that the perspective of $f(u) = u^2$ is $g(u, t) = u^2 / t$ (for $t > 0$),
        and confirm $g$ is convex.
  \item Show that the perspective of $f(u) = -\log u$ (for $u > 0$) is
        $g(u, t) = t \log(t/u)$, and relate this to the KL divergence.
\end{enumerate}
\end{exercise}

\begin{exercise}[Quasiconvex but Not Convex]
\label{ex:quasiconvex-not-convex}
Show that $f(x) = \sqrt{|x|}$ is quasiconvex on $\R$ but not convex.
\textit{Hint}: For quasiconvexity, show sublevel sets are intervals. For non-convexity,
find a counterexample to the convexity inequality.
\end{exercise}

\begin{exercise}[Separating Hyperplane Application]
\label{ex:separation-app}
Let $C = \{x \in \R^2 : x_1^2 + x_2^2 \le 1\}$ (the unit disk) and $p = (2, 0)$.
\begin{enumerate}[label=(\alph*)]
  \item Find a separating hyperplane between $C$ and $\{p\}$.
  \item Find a supporting hyperplane to $C$ at the point $(1, 0)$.
  \item Is the supporting hyperplane unique at this point? Prove your answer.
\end{enumerate}
\end{exercise}

\begin{exercise}[Convexity Preservation Under Partial Minimization]
\label{ex:partial-min-exercise}
Let $f(x, y) = x^2 + y^2 + xy$ defined on $\R^2$.
\begin{enumerate}[label=(\alph*)]
  \item Show that $f$ is convex by computing the Hessian.
  \item Compute $g(x) = \inf_y f(x, y)$ explicitly.
  \item Verify directly that $g$ is convex.
\end{enumerate}
\end{exercise}

%% ========================================================================
%% SECTION 9: SUMMARY TABLE
%% ========================================================================

\section{Summary}

\begin{table}[h!]
\centering
\caption{Summary of key concepts in convex analysis}
\label{tab:summary}
\renewcommand{\arraystretch}{1.3}
\begin{tabular}{@{}p{3.5cm}p{5.5cm}p{5cm}@{}}
\toprule
\textbf{Concept} & \textbf{Definition / Condition} & \textbf{Key Property} \\
\midrule
Convex set &
$\theta x + (1-\theta)y \in C$ for $\theta \in [0,1]$ &
Closed under intersection, affine maps \\

Convex function &
$f(\theta x + (1-\theta)y) \le \theta f(x) + (1-\theta) f(y)$ &
Epigraph is convex set \\

First-order cond. &
$f(y) \ge f(x) + \nabla f(x)^T(y-x)$ &
Local min = global min \\

Second-order cond. &
$\nabla^2 f(x) \succeq 0$ &
Hessian PSD everywhere \\

Jensen's inequality &
$f(\E[X]) \le \E[f(X)]$ for convex $f$ &
Foundation of variational methods \\

Pointwise supremum &
$\sup_\alpha f_\alpha(x)$, each $f_\alpha$ convex &
Preserves convexity \\

Partial minimization &
$\inf_y f(x,y)$, $f$ jointly convex &
Preserves convexity \\

Fenchel conjugate &
$f^*(y) = \sup_x (y^T x - f(x))$ &
Always convex; $f^{**} = f$ if $f$ closed convex \\

Fenchel's inequality &
$f(x) + f^*(y) \ge x^T y$ &
Equality iff $y \in \partial f(x)$ \\

Quasiconvex &
$f(\theta x + (1-\theta)y) \le \max\{f(x), f(y)\}$ &
Sublevel sets convex \\

Log-concave &
$\log f$ concave, $f > 0$ &
Closed under products, marginalization \\

Log-convex &
$\log f$ convex, $f > 0$ &
Implies convexity \\

Separating hyperplane &
$\exists a \ne 0: a^T x \le b \le a^T y$ &
Holds for disjoint convex sets \\

Supporting hyperplane &
$a^T x \le a^T x_0$ for all $x \in C$ &
Exists at every boundary point \\
\bottomrule
\end{tabular}
\end{table}

\vspace{1cm}

\begin{center}
  \textit{End of Lesson 8: Convex Analysis}
\end{center}


\end{document}